% !TEX spellcheck = en_US

% ~~~~~~~~~~~~~~~~~~~~~~~~~~~~~~~~~~~~~~~ V
% (NC) 2026 Haluk Bingol
% github.com/halukbingol/LaTeX-Templates
%
% Licensees may copy, distribute, display, and perform the work and make derivative works
% and remixes based on it only for non-commercial purposes.
% ~~~~~~~~~~~~~~~~~~~~~~~~~~~~~~~~~~~~~~~ A


\documentclass[11pt]{amsart}




% ~~~~~~~~~~~~~~~~~~~~~~~~~~~~~~~~~~~~~~~ V lib
% Note that hblibCommon is loaded by hblibArticle
\usepackage{../cis000_latex/lib/hblibArticle}
%	\usepackage{../cis000_latex/lib/hblibDatastructures}
%	\usepackage{../cis000_latex/lib/hblibCodeoutput}
% ~~~~~~~~~~~~~~~~~~~~~~~~~~~~~~~~~~~~~~~ A lib




% ~~~~~~~~~~~~~~~~~~~~~~~~~~~~~~~~~~~~~~~ V specific
% ~~~~~~~~~~~~~~~~~~~~~~~~~~~~~~~~~~~~~~~ A specific




% ~~~~~~~~~~~~~~~~~~~~~~~~~~~~~~~~~~~~~~~ V specific
% ~~~~~~~~~~~~~~~~~~~~~~~~~~~~~~~~~~~~~~~ A specific




% ~~~~~~~~~~~~~~~~~~~~~~~~~~~~~~~~~~~~~~~ A titleParam
\newcommand{\cisno}{111}	% 111		
\newcommand{\semesterID}{2036-3}	% 2026-3	
\newcommand{\examID}{1m1}	% 1m1	
\newcommand{\questionID}{Q1}	% Q1	
\newcommand{\groupNo}{g1}	% g1	
\newcommand{\topic}{Array Library}	% Array Library
\newcommand{\topicshort}{ArrayLib}		% ArrayLib
\newcommand{\myauthor}{Haluk Bingol}		% Haluk Bingol
\newcommand{\myauthorshort}{Bingol}		% Bingol
% ~~~~~~~~~~~~~~~~~~~~~~~~~~~~~~~~~~~~~~~ A titleParam



% ~~~~~~~~~~~~~~~~~~~~~~~~~~~~~~~~~~~~~~~ V derived
% do not change
\newcommand{\cisxxx}{\textsc{cis\cisno}}	% do not change
% ~~~~~~~~~~~~~~~~~~~~~~~~~~~~~~~~~~~~~~~ A derived



% ~~~~~~~~~~~~~~~~~~~~~~~~~~~~~~~~~~~~~~~ V title
% Common title format for exam
% !TEX spellcheck = en_US

% ~~~~~~~~~~~~~~~~~~~~~~~~~~~~~~~~~~~~~~~ V title
% Common title format for exam question

\title[\courseID-\examNo: \topicShort]{
% ~~~~~~~~~~~~~~~~~~~~~~~~~~~~~~~~~~~~~~~ V box
	\rule[1ex]{\textwidth}{.1pt}
	\begin{tabular*}{\textwidth}{l @{\extracolsep{\fill}} r}
		\courseID\ \courseName
			&\semesterName\ \semesterID\\
		\courseInstructors
			&\examDate\\
		\examName
			&\questionID\\
	\end{tabular*}
	\rule[1ex]{\textwidth}{.1pt}
	~\\
	\vspace{1cm}
% ~~~~~~~~~~~~~~~~~~~~~~~~~~~~~~~~~~~~~~~ A box
	{\Large 
		\questionNo\ . \topic
	}
}

\author{
	\questionAuthor\\
	\hbTimeStamp
}

\begin{document}

\maketitle

\tableofcontents
% ~~~~~~~~~~~~~~~~~~~~~~~~~~~~~~~~~~~~~~~ A title

% ~~~~~~~~~~~~~~~~~~~~~~~~~~~~~~~~~~~~~~~ A title




% ~~~~~~~~~~~~~~~~~~~~~~~~~~~~~~~~~~~~~~~ V
\textbf{Warning.} Please read \refapp{sec:appExaminfo} Exam Info.
% ~~~~~~~~~~~~~~~~~~~~~~~~~~~~~~~~~~~~~~~ A


% ~~~~~~~~~~~~~~~~~~~~~~~~~~~~~~~~~~~~~~~ V
\section{Introduction}
\label{sec:intro}

Some intro



% ~~~~~~~~~~~~~~~~~~~~~~~~~~~~~~~~~~~~~~~ V
\hSectionQuestion



% ~~~~~~~~~~~~~~~~~~~~~~~~~~~~~~~~~~~~~~~ V Q1
\subsection{\hCode{sumOfNeighbors} method}
\label{sec:sumOfNeighbors}



% ~~~~~~~~~~~~~~~~~~~~~~~~~~~~~~~~~~~~~~~
\subsubsection{Signature}

% +++++++++++++++++++++++++++++++++++++++ V
%: -Code.
\begin{javacode}
public static int[] sumOfNeighbors(int[] arr)
\end{javacode}
% +++++++++++++++++++++++++++++++++++++++ A



% ~~~~~~~~~~~~~~~~~~~~~~~~~~~~~~~~~~~~~~~
\subsubsection{Function} 

Description of the function.



% ~~~~~~~~~~~~~~~~~~~~~~~~~~~~~~~~~~~~~~~
\subsubsection{Rules}

Some rules.



% ~~~~~~~~~~~~~~~~~~~~~~~~~~~~~~~~~~~~~~~
\subsubsection{Example}

Some examples.



% ~~~~~~~~~~~~~~~~~~~~~~~~~~~~~~~~~~~~~~~
\subsubsection{Constraints}

Some constraints.

\textbf{Hint.} 
Some hint.


% ~~~~~~~~~~~~~~~~~~~~~~~~~~~~~~~~~~~~~~~ A Q1




% ~~~~~~~~~~~~~~~~~~~~~~~~~~~~~~~~~~~~~~~ V
\hNumbersArabic % back to numbers
% ~~~~~~~~~~~~~~~~~~~~~~~~~~~~~~~~~~~~~~~ A


% ~~~~~~~~~~~~~~~~~~~~~~~~~~~~~~~~~~~~~~~ V testing
\section{Testing}

Use \hCode{TS\_xxx\_jUnit} to test your code.
You can also use \hCode{xxx\_Test} if there is any.
% ~~~~~~~~~~~~~~~~~~~~~~~~~~~~~~~~~~~~~~~ A testing




% ~~~~~~~~~~~~~~~~~~~~~~~~~~~~~~~~~~~~~~~ V
\section{ToDo}\label{sec:todo}

\begin{enumerate}
	\item Fill in your personal information at the beginning of \hCode{ArrayLib.java}.
	\item Complete the \hCode{sumOfNeighbors} method.
	\item Complete the \hCode{productOfOdds} method.
	\item Complete the \hCode{selectAtMost} method.
	\item Write your code only between the \hCode{below} and \hCode{above} marks.
	\item Run \hCode{TS\_ArrayLib\_jUnit} to verify your implementation.
\end{enumerate}
% ~~~~~~~~~~~~~~~~~~~~~~~~~~~~~~~~~~~~~~~ A




% ~~~~~~~~~~~~~~~~~~~~~~~~~~~~~~~~~~~~~~~ V \appendix
%\section*{At the end}\label{sec:end}

%See \hCode{cis11x\_ExamInfo}.
% Exam info

\newpage
\appendix


% ~~~~~~~~~~~~~~~~~~~~~~~~~~~~~~~~~~~~~~~ V
\section{Exam}
% ~~~~~~~~~~~~~~~~~~~~~~~~~~~~~~~~~~~~~~
\subsection{Introduction}

An exam is composed of three periods:
\begin{enumerate}[label=\roman*.]

	\item \textbf{Period-1.} 
	The internet is on. You should complete all downloading in Period-1.

	\item \textbf{Period-2.} 
	The internet is off. You should be solving the questions.

	\item \textbf{Period-3. }
	The internet is on. You should zip your solution and upload to \YULearn.
	
\end{enumerate}
% ~~~~~~~~~~~~~~~~~~~~~~~~~~~~~~~~~~~~~~~ V
\subsection{Period-1. Before exam starts, (internet on)}

The internet is on for downloading.

Start with a clean system. That is,

\begin{enumerate}[label=\roman*.]
	
	\item \textbf{Communication devices.}
	Leave all communication devices, including phones, smart watches,
	tablets, on the instructor desk.

	\item \textbf{Restart.} 
	Shut down and restart your computer if it is not already restarted.
	
	\item \textbf{eclipse -clean.} 
	Sometimes eclipse needs a clean start. 
	Open terminal and run the following

% +++++++++++++++++++++++++++++++++++++++ V lst
%: -Lst.
\begin{lstlisting}[style=codesnippet, language=bash]
C:\app\eclipse-java-web\eclipse.exe -clean
\end{lstlisting}
% +++++++++++++++++++++++++++++++++++++++ A lst

	
	\item \textbf{Create a new workspace.} 
	Create a brand-new eclipse workspace
	\begin{enumerate}[itemsep=1pt, topsep=1pt] 
		
		\item \menu{File > Switch Workspace > Other}
	
		\item Click \menu{Browse}, go to \menu{Desktop}
	
		\item Click \menu{New Folder} to create a new folder, name it as \hCode{Snnn}, 
		where $nnn$ is your student no.
		
	\end{enumerate}

	
	\item \textbf{Create a new Java project.} 
	Create a brand-new Java project.
	\begin{enumerate}[itemsep=1pt, topsep=1pt] 
		
		\item \menu{File > New > Other....}
	
		\item Select \menu{Java Project}.
	
		\item Set Project name as your name or \hCode{Sxxx}, where xxx is your student no.
	
		\item Make sure that \menu{Create module-info.java} file is not selected. 
		
	\end{enumerate}
	
		
	Hence, there is nothing under \hCode{src} at the beginning. 
	This becomes important when you zip
	\hCode{src} for submission at the end of the exam.
	
	\item \textbf{Add jUnit.} 
	Add \hCode{jUnit} 5 library to your project.
	
	\item \textbf{question.zip.}
	For simplicity, let \hCode{question.zip} be the exam question. 
	Actually, the downloaded exam question has a file name 
	like \hCode{cis112\_2025\_e1M1\_Q1.zip}.
	\begin{enumerate}[itemsep=1pt, topsep=1pt] 
		
		\item Download \hCode{question.zip} package. 
		Since it is password protected, 
		you need to wait for the password.
		
	\end{enumerate}
	
	\emph{
		Make sure that you downloaded \hCode{question.zip} 
		before the internet is disconnected.
	}

\end{enumerate}
% ~~~~~~~~~~~~~~~~~~~~~~~~~~~~~~~~~~~~~~~ A
% ~~~~~~~~~~~~~~~~~~~~~~~~~~~~~~~~~~~~~~~ V
\subsection{Period-2. During exam (internet off)}

The exam starts after the internet is disconnected.

\begin{enumerate}[label=\roman*.]  
	
	\item \textbf{Import question.}
	Note that \hCode{question.zip} is password protected. 
	After the password is shared, you can import it in eclipse. 
	After the required packages are imported, 
	your Package Explorer in eclipse should be like \reffig{fig:DirectoryStructure}
	





% +++++++++++++++++++++++++++++++++++++++ V
%: -Fig.
\begin{figure}[htbp]
\begin{center}
	\begin{forest}
%	for tree={font=\ttfamily, grow'=0, folder, draw=none}
	for tree={font=\ttfamily, grow'=0, folder, draw=none,
	          s sep=0pt,        % no extra gap between lines
	          inner ysep=1pt}   % less padding around each name
		[project/
			[src/                     <------------------- zip at this level
				[mypackage/ 
					[question.pdf] 
					[ModifyMe.java] 
					[DoNotModifyMe.java] 
					[TS\_question\_jUnit.java]
				]
				[ppp/ 
					[Aaa.java] 
					[Bbb.java]
					[kkk/
						[Ccc.java]
					] 
				]
			]
		]
	\end{forest}

	\caption{\small
		\textbf{Directory structure.}
		The codes are under package \hCode{mypackage}.
		It may be the case that there are multiple packages 
		such as  \hCode{mypackage}, \hCode{ppp} and \hCode{ppp/kkk}.
		Note that class \hCode{ModifyMe.java} has 
%\begin{javacode}
%package ppp.kkk;}
%\end{javacode}		
		``\hCodeJ{package mypackage;}''
		as the first line;
		while class \hCode{Ccc} has 
		``\hCodeJ{package ppp.kkk;}''.
		%
		\hCode{question.pdf} is the description of the question.
		\hCode{ModifyMe.java} is the file you are expected to modify.
		There may be multiple such files.
		\hCode{DoNotModifyMe.java} is the file that is used but 
		you are not allow to modify.
		There may be multiple such files.
		\hCode{TS\_question\_jUnit.java} is the jUnit file shared with you 
		so that you can test your development.
	}
	\label{fig:DirectoryStructure}
\end{center}
\end{figure}
% +++++++++++++++++++++++++++++++++++++++ A
	
	\begin{enumerate}[itemsep=1pt, topsep=1pt] 
		
		\item The password of \hCode{question.zip} is given.
		
		\item Uncompress \hCode{question.zip} using the given password.
		
		\item Import \hCode{question.zip} package(s) into eclipse. 
		Make sure that the package is under \hCode{src}.
		\item Start reading the description pdf.
		
	\end{enumerate}

	\item \textbf{Open pdf files.} 
	It is not recommended to opening \hCode{.pdf} files within eclipse.
	\begin{enumerate}
		
		\item Right click on the \hCode{.pdf} file.
		
		\item \menu{Open With > Default Editor} (or \menu{System Editor})
		
	\end{enumerate}

\end{enumerate}
% ~~~~~~~~~~~~~~~~~~~~~~~~~~~~~~~~~~~~~~~ A




% ~~~~~~~~~~~~~~~~~~~~~~~~~~~~~~~~~~~~~~
\subsubsection{Student info}
\label{sec:app:studentInfo} 

Make sure that you populate \hCode{Student info} in \hCode{modifyMe.java} files.
Be careful not to change ``below'' and ``above'' marks.

% +++++++++++++++++++++++++++++++++++++++ V
%: -Code.
\input{../../cis000_latex/input/in-CodeYourInfo.tex}
% +++++++++++++++++++++++++++++++++++++++ A




% ~~~~~~~~~~~~~~~~~~~~~~~~~~~~~~~~~~~~~~
\subsubsection{Where to change: ``below-above'' marks} 
\label{sec:app:belowAboveMarker}

You are only allowed to change between ``below'' and ``above'' marks in \hCode{modifyMe.java} files.
\emph{Be careful not to change} ``below'' and ``above'' marks.

% +++++++++++++++++++++++++++++++++++++++ V
%: -Code.
\input{../../cis000_latex/input/in-CodeBelowAboveMarker}
% +++++++++++++++++++++++++++++++++++++++ A

% ~~~~~~~~~~~~~~~~~~~~~~~~~~~~~~~~~~~~~~~ V
\subsection{Period-3. At the end of the Exam (internet on)}

The internet is on for uploading.

\subsection{Uploading}

\begin{enumerate}[label=\roman*.]
	
	\item Right click on src.

	\item Select \menu{Show in > System Explorer} to open 
	\hCode{src} in the \hCode{System Explorer}.

	\item Compress \hCode{src} directory in zip format. 
	Make sure that the name of the compressed file is \hCode{src.zip}.

	\textbf{Warning.}
	\begin{itemize}
		
		\item Only \hCode{zip} compression is accepted. 
		Other compressions, such as \hCode{rar}, are not accepted.

		\item Do not zip
		\begin{itemize}[nosep] 

			\item workspace

			\item project

			\item package
			
		\end{itemize}
			
		\end{itemize}
		but \hCode{src} only.

	\item Upload \hCode{src.zip} to \YULearn.

	\item Zip \hCode{src} directory of eclipse and obtain \hCode{src.zip}.

	\item Check if \hCode{src.zip} is correct
	\begin{enumerate}
		
		\item Create a new folder on Desktop
		
		\item Copy \hCode{src.zip} to the new folder
		
		\item Expand \hCode{src.zip}. 
		The directory structure given in \reffig{fig:DirectoryStructure} should be obtained.
		Otherwise AutoGrader will assign 0, 
		since it cannot locate your files to grade.
	\end{enumerate}

	\item Upload \hCode{src.zip} to YULearn.
		
\end{enumerate}
% ~~~~~~~~~~~~~~~~~~~~~~~~~~~~~~~~~~~~~~~ A


\subsection{Warning about the AutoGrader}
\label{sec:app:warnAutoGrader}

Note that the AutoGrader will assign a 0 grade 
if uploaded zip 

\begin{itemize}[nosep] 
	
	\item does not expand as given in \reffig{fig:DirectoryStructure}.
	
	\item has compile-time errors.
	
	\item has infinite loops.
	
\end{itemize}


% ~~~~~~~~~~~~~~~~~~~~~~~~~~~~~~~~~~~~~~ V
\subsection{Timing}


%\begin{table}[htp]
%\caption{default}
\begin{center}
	\begin{tabular}{|l|r|l|}

\hline
\textbf{Time} & \textbf{min} & \textbf{Description} \\
\hline

t + 00:00  & 15 & \textbf{Period-1. Getting ready} \\
 &  & - Leaving phones on the instructor desk \\
 &  & - Student seating \\
 &  & - Restart of the computers, \\
 &  & - Downloading question zip from YULearn \\
 
t + 00:15  & 75 & \textbf{Period-2. Exam begins} \\
 &  & - The Internet is disconnected \\
 &  & - Password of question zip is shared \\
 &  & - If you finished early, prepare your \hCode{src.zip}. \\
t + 01:30  & 10 & Exam ends \\
 &  & - Students collected next to the board \\
 &  & - jUnit grading \\
 
t + 01:40  & 10 & \textbf{Period-3. Finalize}  \\
 &  & - The Internet is reconnected \\
 &  & - \hCode{src.zip} is uploaded to YULearn \\
t + 01:50  & 10 & Leave the room \\
t + 02:00  &  & Start of the next exam \\

\hline

 & \textbf{120} & \textbf{Total} \\
 
\hline

 	\end{tabular}
\end{center}
%\label{default}
%\end{table}%
% ~~~~~~~~~~~~~~~~~~~~~~~~~~~~~~~~~~~~~~ A
	
% ~~~~~~~~~~~~~~~~~~~~~~~~~~~~~~~~~~~~~~~ V
\section{Objection}

Prepare to \textbf{loose points} 
if your objection ends with no increase in your grade. 
In order to avoid this,
check your submission with the following procedure.




% ~~~~~~~~~~~~~~~~~~~~~~~~~~~~~~~~~~~~~~
\subsection{Test your submission}

Suppose \hCode{doit.java} is the file student should complete during exam. 
If there are more than one such files, repeat the procedure for each. 
Only then run the tests.
Note that your \hCode{doit.java} is tested in the environment of \hCode{question}. 
That is,
\begin{itemize}[nosep] 
	
	\item all the files comes from \hCode{question}.
	
	\item Only \hCode{doit.java} files come from your submission.
	
\end{itemize}
If you modify files that are not in the list of \hCode{doit.java} files, 
these changes no contribution to the grading 
since only original question files are used. 
Note that if you do this, 
you code may not compile.

Hence, \hCode{TS\_ file} in your submission is not used.




% ~~~~~~~~~~~~~~~~~~~~~~~~~~~~~~~~~~~~~~
\subsubsection{Get question}

\begin{enumerate}[label=\roman*.]
	
	\item Download the \hCode{question.zip} from \YULearn.
	\item Expand \hCode{question.zip} to obtain question.
	\item In eclipse, create a brand new java project \hCode{ProjObjection}.
	\item Import question into \hCode{ProjObjection}.
	\item Select \hCode{doit.java}, copy-paste to obtain doit\_\_question.java
	
\end{enumerate}




% ~~~~~~~~~~~~~~~~~~~~~~~~~~~~~~~~~~~~~~
\subsubsection{Get submission}
\begin{enumerate}
	
	\item Download your submission \hCode{src.zip} from \YULearn.
	\item Expand \hCode{src.zip} to obtain submission.
	
	\item Select \hCode{doit.java} in submission. 
	Copy-paste it in eclipse. 
	Note that \hCode{doit.java} of submission will overwrite \hCode{doit.java} of question. 
	Eclipse will warn you.
	
\end{enumerate}




% ~~~~~~~~~~~~~~~~~~~~~~~~~~~~~~~~~~~~~~
\subsubsection{Testing}

Now, setup is ready to be tested.

\begin{enumerate}
	
	\item Test \hCode{doit.java} using the corresponding hidden test file \hCode{TH\_ }.

	\item Check the number of tests passed. 
	If this number is different that your grade, 
	initiate objection by an email containing
	\begin{itemize}[nosep] 
		
		\item your name, last name, student number,
		\item courseCode, section,
		\item exam and question no that you are objecting,
		\item your grade and what you get at the end of this procedure,
		\item explanation the situation.
	\end{itemize}

	
\end{enumerate}

% ~~~~~~~~~~~~~~~~~~~~~~~~~~~~~~~~~~~~~~~ A

% ~~~~~~~~~~~~~~~~~~~~~~~~~~~~~~~~~~~~~~~ A

% ~~~~~~~~~~~~~~~~~~~~~~~~~~~~~~~~~~~~~~~ A \appendix



\end{document}
